% !TeX root = ./main.tex

\begin{abstract}

传统文本检索增强生成（RAG）在处理图文混排、跨页引用和结构复杂的长文档时，容易丢失图表、版面和跨页关系等关键信息，导致答案缺乏稳定证据支撑。针对这一问题，本文设计并实现了 DocRAG————一个多模态 GraphRAG 文档问答系统，采用``FAISS 向量索引 + networkx 异构文档图 + 视觉语言模型''的技术路线，完成从 PDF 入库、页面结构化解析、向量索引与知识图谱构建到多模态检索问答的完整链路。系统使用 PyMuPDF 渲染页面图像，调用 qwen3-vl-8b-instruct 抽取文本、表格、图表、结论、关键词和关系，并在问答阶段融合向量召回、图谱邻域扩展和页面图像证据生成答案。

实验方面，本文基于 pdfQA 中的 12 篇 ESG 报告构建了 120 道覆盖文本定位、表格问答、图表理解、跨页延续、多跳关系和引用追踪的问题集，并设计五组消融实验评估各组件作用。结果表明，完整系统在证据页命中和答案生成稳定性上表现最好，FAISS 语义召回是核心基础，知识图谱和页面图像进一步提升了证据定位与视觉推理能力。Token 成本实验显示，DocRAG 仅发送检索后的相关页面和文本证据，可将单次问答输入从约 5--8 万等效 token 降至约 3000--5000 等效 token，输入 token 节省约 94\%，显著降低了长文档多轮问答和批量评测的调用成本。

\end{abstract}
